\begin{theorem}[Wild quadratic nonboundary cancellation]\label{mod:cases-wildquadratic-nonboundary}
Let \(K/F\) be the prime-cyclic quadratic local extension in residue
characteristic \(2\), with lower break \(t\), residue degree one, and a
chosen integral uniformizer \(\pi_K\) which generates
\(\mathcal O_K\) over \(\mathcal O_F\).  Put

\[
  T=t+1,
\]

and let \(m>1\) be the conductor indexing a coherent
wild-quadratic common coefficient view.  The theorem has two entry points.

For an exact \texttt{WildQuadraticAllEvenCoefficientView}, supply the
computational data, phase data, real
\texttt{ExactQuadraticPhaseAssembly}, the four actual even Lamprecht rows,
the q-free common coordinate compatibility, the existing
\texttt{CompleteCorrectionData}, the finite Delta hypotheses, and
whole-orbit minimality of the base character.  The theorem
\texttt{wildQuadratic\_allEven\_elementary} reuses the actual
\(\tau\)-, base-, and twist-row conductor equalities, the order of \(A\),
the trace-pullback additive conductor, and the correction-depth argument.
At the even boundary it derives the order of \(1+n\) from the same actual
selected \(\tau\)- and \(\tau\chi_F\)-pairs and minimality; no new
denominator-order hypothesis is introduced.  Combining this elementary
result with \texttt{wildQuadratic\_allEvenErrorFormula} gives
\texttt{wildQuadratic\_nonboundary\_allEven}:
\[
 \operatorname{Err}_{K/F}(\chi_F)=1.
\]
This entry point contains no normalized refinement and applies exactly to
\[
 \texttt{belowMEvenTEven},\qquad
 \texttt{boundaryEven},\qquad
 \texttt{aboveMEvenTEven}.
\]

For the complementary compatibility interface, let
\texttt{WildQuadraticCoefficientView} \(V\) be the existing refined view.
Supply its literal phase-row and coordinate compatibilities, refinement
assembly, finite Delta hypotheses, and whole-orbit minimality.  Assume that
its row is not the exceptional odd boundary row.  The existing public
\texttt{wildQuadratic\_nonboundary} theorem and signature are unchanged.
In either entry point,

\[
 \operatorname{Err}_{K/F}(\chi_F)=1.
\]

For the refined entry point, the nonboundary condition is proved to mean
exactly one of the following nine rows, with no endpoint absorbed into a
strict range:

\[
\begin{array}{llll}
 \texttt{belowMEvenTEven},&\texttt{belowMEvenTOdd},&
 \texttt{belowMOddTEven},&\texttt{belowMOddTOdd},\\
 \texttt{boundaryEven},\\
 \texttt{aboveMEvenTEven},&\texttt{aboveMEvenTOdd},&
 \texttt{aboveMOddTEven},&\texttt{aboveMOddTOdd}.
\end{array}
\]

Equivalently, their signed-depth alternatives are \(m<T\), \(m=T\) with
\(T\) even, and \(T<m\).  The row \(m=T\) with \(T\) odd is excluded and no
theorem about that boundary is used.

For the depth calculation put \(a=T-m\).  The common correction is retained
literally as

\[
  X=\frac{2s}{1+n},\qquad s=\operatorname{Tr}_{K/F}(u),
  \qquad v_K(u)=v_F(n)=a.
\]

In equal characteristic \(2\), \(X=0\).  Otherwise the trace-ideal theorem
gives the exact bounds

\[
  2\in\mathfrak p_F^{\lfloor T/2\rfloor},\qquad
  s\in\mathfrak p_F^{\lfloor(a+T)/2\rfloor}.
\]

The proof keeps the denominator \(1+n\) until its order is justified.  Its
order is \(0\) when \(m<T\), is \(a\) when \(T<m\), and is \(0\) at the even
boundary.  The last assertion is not assumed: it is derived from the two
actual selected stationary ratios for the \(\tau\)- and
\(\tau\chi_F\)-rows, the exact conductor formula
\(m(\tau\chi_F)=\max(m,T)\), and whole-orbit minimality.  The corresponding
floor inequalities, including the parity equality at the even boundary,
then prove

\[
  X\in\mathfrak p_F^T.
\]

The selected numerator in every supplied Lamprecht row has order zero, so
its actual selected ratio has order \(-(m(\theta)+n(\psi_F))\).  Applied to
the real \(\tau\)-row this gives, with the same literal representative used
by the exact error formula,

\[
  v_F(A)=-(T+n_F(\psi_F)).
\]

No canonical stationary representative is chosen.  The pullback-conductor
theorem is also instantiated at this exact depth and records

\[
 n_K(\psi_F\circ\operatorname{Tr}_{K/F})
   =2n_F(\psi_F)+T.
\]

Consequently the conductor triviality lattice kills the retained elementary
phase in its actual orientation:

\[
  \psi_F(-AX)=1.
\]

Finally, the exact error formula evaluates all nine finite nonboundary
coefficient rows.  With the positive refinement/Hasse orientation used
there, it proves

\[
  \Phi_{\mathrm{table}}R_\chi R_\tau=1,
  \qquad
  \operatorname{Err}_{K/F}(\chi_F)=\psi_F(-AX).
\]

Thus the table phases and the two refinement/Hasse corrections cancel
together, while the separately retained elementary phase is one, and the
error term is exactly one.  Neither the boundary theorem nor the assembled
wild-quadratic main theorem is assumed.
\LeanFile{LanglandsFirstMainLemma/Cases/WildQuadratic/NonBoundary.lean}
\lean{LanglandsFirstMainLemma.wildQuadratic_nonboundary_rows}
\lean{LanglandsFirstMainLemma.wildQuadratic_nonboundary_depthCase}
\lean{LanglandsFirstMainLemma.LocalLamprechtPhaseData.selectedStationaryPair_beta_order}
\lean{LanglandsFirstMainLemma.LocalLamprechtPhaseData.selectedStationaryPair_ratio_order}
\lean{LanglandsFirstMainLemma.wildQuadratic_correction_mem}
\lean{LanglandsFirstMainLemma.wildQuadratic_boundary_denominator_order}
\lean{LanglandsFirstMainLemma.wildQuadratic_A_order}
\lean{LanglandsFirstMainLemma.wildQuadratic_additiveConductor_compTrace}
\lean{LanglandsFirstMainLemma.wildQuadratic_allEven_elementary}
\lean{LanglandsFirstMainLemma.wildQuadratic_nonboundary_elementary}
\lean{LanglandsFirstMainLemma.wildQuadratic_nonboundary_allEven}
\lean{LanglandsFirstMainLemma.wildQuadratic_nonboundary}
\leanok
\SourceText{Lemma lem:quadratic-depth and Proposition prop:wild-quadratic-cancellation, nonboundary cases.}
\uses{mod:cases-wildquadratic-errorformula,mod:ramification-pullbackconductors}
\FormalizationContract
The q-free theorem is restricted to the exact three all-even rows, and the
existing refined theorem remains restricted to the exact nine rows listed
above.  The inequalities \(m<T\) and \(T<m\) remain strict, while equality
is admitted only with an explicit proof that \(T\) is even.  Both paths
derive the orders of \(2\), the trace \(s\), the denominator \(1+n\), and
the actual stationary ratio at the depths shown; neither replaces them by
weaker conductor bounds.

The quadratic unit correction is preserved through the denominator-order
argument.  The exact additive trace-pullback conductor is obtained from
module~\ref{mod:ramification-pullbackconductors}.  The final phase has the
minus orientation \(\psi_F(-AX)\).  Stationary pairs are the supplied
representatives of quotient classes and are compared literally; no
canonical representative or representative-independent isolated phase is
introduced.  The final equality is obtained by combining the exact
nonboundary contract of module~\ref{mod:cases-wildquadratic-errorformula}
with the proved conductor triviality, without importing or invoking the
odd-boundary or wild-quadratic main theorem.
\end{theorem}
